\section{Nonzero endpoint trace on every eigenspace}
\label{sec:spectral-trace-multiplicity}

\begin{theorem}[AC-S44: the trace kernel is exactly the Sobolev part]
\label{AC-S44}
Fix an arbitrary window and an eigenvalue \(\alpha\) of its
actual real Weil operator \(A\).  Write
\(E_\alpha=\ker(A-\alpha)\) for the real eigenspace and
\(m=\dim E_\alpha\).  Then the endpoint trace on \(E_\alpha\)
has rank exactly one.  More precisely, there is a sign
\(\epsilon_\alpha\in\{1,-1\}\) such that
\[
 \tau_-=\epsilon_\alpha\tau_+\quad\hbox{on }E_\alpha,
 \qquad \tau_+|_{E_\alpha}\ne0.                         \tag{TM.1}
\]
Its common kernel is
\[
 K_\alpha=\ker(\tau_+|_{E_\alpha})
          =E_\alpha\cap H^1_0(I),\qquad
 \dim K_\alpha=m-1,
 \qquad \partial_xK_\alpha\subset E_\alpha.              \tag{TM.2}
\]
In particular every eigenvalue, including the lowest one on
every window, has an eigenvector with nonzero endpoint trace.
Every simple eigenvalue has nonzero trace on every nonzero
eigenvector.  Simplicity is equivalently expressed by
\(E_\alpha\cap H^1_0(I)=\{0\}\); this equivalence does not
establish that exclusion for the general lowest eigenspace.
All assertions are unconditional.
\end{theorem}
\begin{proof}
\proofstep{1. A bounded-forcing reduced resolvent.}
Let \(P_0\) project onto \(E_\alpha\), and let
\(R_\alpha\) be \((A-\alpha)^{-1}\) on \(E_\alpha^\perp\).
For bounded \(g\perp E_\alpha\),
\[
       R_\alpha g\in\mathcal X_I,
 \qquad
 \|R_\alpha g\|_\infty+
            \|G_I R_\alpha g\|_\infty
                         \leq C_\alpha\|g\|_\infty.     \tag{TM.3}
\]
Here the constant may depend on the fixed window and eigenvalue.
To prove this, let \(P_-\) and \(P_+\) be the spectral
projections strictly below and strictly above \(\alpha\).
The below-eigenvalue part
\[
      \sum_{\alpha_j<\alpha}
              \frac{\langle u_j,g\rangle}{\alpha_j-\alpha}u_j
\]
is a finite sum of bounded eigenfunctions, by AC-S36.
For the other part use
\[
 R_\alpha P_+g
       =\int_0^\infty e^{-t(A-\alpha)}P_+g\,dt.          \tag{TM.4}
\]
Define the integral first in \(L^2\).  The bounds below also
make its truncated integrals Cauchy in \(L^\infty\), by testing
against \(L^1\cap L^2\); no strong continuity of the semigroup
on all of \(L^\infty\) is needed.
For small \(t\), the killed positive semigroup of \(G_I\)
and the bounded \(L^\infty\) perturbation in (RN.4) give
\(\|e^{-tA}\|_{\infty\to\infty}\leq e^{Ct}\).
Subtracting the finite bounded spectral projections
\(P_-+P_0\) gives a uniform small-time bound for the
integrand in (TM.4).
For large \(t\), choose \(t_0>1/2\).  AC-S36 gives
\(e^{-t_0A}:L^2\to L^\infty\), and the strictly positive
gap above \(\alpha\) gives exponential decay of
\(e^{-(t-t_0)(A-\alpha)}P_+\) on \(L^2\).
These estimates prove the asserted \(L^\infty\) convergence.
The spectral theorem puts \(w=R_\alpha g\) in
\(\operatorname{Dom}(A)\), with \(Aw=\alpha w+g\).
Since \(A=G_I+B\) with \(B\) bounded on \(L^\infty\),
this proves (TM.3).

\proofstep{2. Zero trace forces an eigenspace derivative.}
Let \(u\in E_\alpha\) have both endpoint traces zero.
For every real \(g\in C_c^\infty(I)\cap E_\alpha^\perp\),
put \(w=R_\alpha g\).  It belongs to \(\mathcal X_I\)
by step~1, so (BT.3) gives
\[
                         \int_Iu g'=0.                  \tag{TM.5}
\]
We justify the distributional annihilator explicitly.
Choose a real orthonormal basis \(v_1,\ldots,v_m\) of
\(E_\alpha\).  By density of \(C_c^\infty(I)\) in \(L^2(I)\),
choose real compact smooth \(\psi_1,\ldots,\psi_m\)
such that the matrix
\(M_{ij}=\langle v_i,\psi_j\rangle\) is invertible;
for example, sufficiently close \(L^2\) approximations to
the basis have this property.
For any compact smooth \(g\), subtract
\(\sum_j c_j\psi_j\), where
\(c=M^{-1}(\langle v_i,g\rangle)_i\).
The result lies in the test space of (TM.5).
Consequently the functional \(g\mapsto\int ug'\)
is a linear combination of \(g\mapsto\langle v_i,g\rangle\).
Thus the distributional derivative \(u'\) belongs to
\(E_\alpha\subset L^2(I)\).
It follows that \(u\in H^1(I)\).
Its continuous representative has zero endpoint values by
AC-S36, hence \(u\in H^1_0(I)\).

Conversely, if \(u\in E_\alpha\cap H^1_0(I)\), then
\[
 |u(d)|\leq\sqrt d\,\|u'\|_2,
 \qquad |u(\ell-d)|\leq\sqrt d\,\|u'\|_2.
\]
Multiplication by \(\sqrt{\log(1/d)}\) shows that both
renewal traces vanish.  We have proved the equality in
(TM.2), as well as the derivative inclusion, before using
nonvanishing of the trace on the whole eigenspace.

\proofstep{3. The entire trace cannot vanish.}
AC-S43-trace-rank says that the joint trace rank is at most one.
If it were zero, step~2 would give
\(\partial_xE_\alpha\subset E_\alpha\).
Complexify this finite-dimensional space.  The induced derivative
matrix has an eigenvector \(0\ne v\in E_\alpha\otimes\mathbb C\)
and an eigenvalue \(c\in\mathbb C\), so
\(v'=cv\) as distributions on \(I\).
The already proved \(H^1_0\) regularity shows that
\(e^{-cx}v(x)\) is an \(H^1\) function with zero derivative.
It is constant and has zero endpoint value, hence is zero.
This contradicts \(v\ne0\).
The trace therefore has rank exactly one.
The fixed-sign conclusion in AC-S43-trace-rank gives (TM.1),
and rank-nullity gives \(\dim K_\alpha=m-1\).
All the remaining assertions follow directly from (TM.2).
\end{proof}

\begin{corollary}[AC-S44-parity: the two parity multiplicities]
\label{AC-S44-parity}
On \(E_\alpha\), let \(p\) be the dimension of the
reflection parity \(\epsilon_\alpha\) in (TM.1), and let
\(o\) be the dimension of the opposite parity.  Then
\[
                              o\leq p\leq o+1.          \tag{TM.6}
\]
The endpoint trace is supported on the parity of dimension \(p\).
This restriction does not decide whether the lowest eigenvalue
is simple or whether a simple lowest eigenvector is even.
\end{corollary}
\begin{proof}
Reflection interchanges \(\tau_+\) and \(\tau_-\), so
(TM.1) makes the trace vanish on the opposite parity and
be nonzero on parity \(\epsilon_\alpha\).
Differentiation reverses parity and is injective on
\(H^1_0(I)\): a zero derivative is a constant with zero
endpoint values.  Apply (TM.2) first to the whole opposite
parity space, and then to the \((p-1)\)-dimensional trace
kernel in parity \(\epsilon_\alpha\).
The resulting injections give \(o\leq p\) and \(p-1\leq o\).
\end{proof}

\paragraph{The remaining lowest-line issue.}
The nonvanishing theorem concerns the endpoint functional on
each eigenspace.  If the eigenspace has dimension greater than
one, it still contains the \((m-1)\)-dimensional zero-trace
space in (TM.2).  Its derivative need not stay in that kernel;
it may enter the trace-carrying direction.  No argument above
excludes this possibility for the lowest eigenspace.
Thus the Sobolev characterization is an equivalent formulation
of simplicity, not a proved new exclusion or a large-window
ground-line selection estimate.
